% BEGIN GENERATED ARTIFACT: prf:conditional-equivalence-laws-propositional-logic
\newpage
\phantomsection
\label{prf:conditional-equivalence-laws-propositional-logic}
\LRAProofFor{thm:conditional-equivalence-laws-propositional-logic}

\begin{remark*}[Return]
\hyperref[thm:conditional-equivalence-laws-propositional-logic]{Return to Theorem (Conditional Equivalence Laws).}
\end{remark*}

\begin{theorem*}[Conditional Equivalence Laws]
For all \(\varphi,\psi\in\WFF\),
\[
\varphi\to\psi
\equiv
\neg\varphi\lor\psi,
\]
and
\[
\neg(\varphi\to\psi)
\equiv
\varphi\land\neg\psi.
\]
\end{theorem*}

\begin{proof}
\textbf{Professional Standard Proof.}~
TODO: supply the compact proof.
\end{proof}

\begin{proof}
\textbf{Detailed Learning Proof.}~
TODO: supply the detailed learning proof.
\end{proof}

\begin{remark*}[Proof structure]
Use the semantic truth clause for implication.
\end{remark*}

\begin{dependencies}
\begin{itemize}
  \item \hyperref[def:logical-equivalence-propositional-logic]{Logical Equivalence}
  \item \hyperref[thm:unique-extension-truth-assignment-propositional-logic]{Unique Extension of a Truth Assignment}
\end{itemize}
\end{dependencies}

\clearpage
% END GENERATED ARTIFACT: prf:conditional-equivalence-laws-propositional-logic
